\documentclass[11pt]{article}
\usepackage[margin=1in]{geometry}
\usepackage{amsmath,amssymb,amsthm,hyperref}
\newtheorem{theorem}{Theorem}
\newtheorem{lemma}[theorem]{Lemma}
\newtheorem{proposition}[theorem]{Proposition}
\newtheorem{definition}[theorem]{Definition}
\newcommand{\SAT}{\operatorname{SAT}}
\title{COST-CERTIFICATES-0001: Charged Closure and a Certificate-Length Boundary}
\author{P versus NP project review}
\date{10 October 2026}
\begin{document}
\maketitle
\paragraph{Status.} This document contains elementary conditional composition
lemmas, an explicit diagonalization obstruction, and open construction targets.
It is not a proof of $P\ne NP$. Mathematical novelty relative to the entire
project or literature is not claimed. The diagonalization boundary below is a
new deduction in this review. REDTEAM independently accepted its proof and
the charged composition after the termination and finite-length conventions
recorded here were made explicit.

\section{Parameters and charging convention}
Let $N$ be the complete encoded bit length of an ordinary CNF, rather than its
variable count. Let $A$ be one fixed exact SAT decider with time $N^{K+o(1)}$
and writable space $N^{\sigma+o(1)}$. Its finite code may be hardwired;
length-dependent advice may not. For each actually executed phase, charge
\[
 W=D+V+U+R+E+X+C,
\]
where $D$ is successful and failed discovery, $V$ verification, $U$ update,
$R$ replay, $E$ emission, $X$ physical source/child access, and $C$ combination.
Terms describe disjoint charged operations, to avoid double counting.
Peak space includes all simultaneously live states, buffers, indexes,
certificates, and stack frames. A reused result is charged once only if its
reuse is implemented; recomputation is charged each time.

\section{Aggregate closure with physical access}
The following restates the audited Recovery7/Recovery10 accounting with the
interface hypotheses made explicit.
\begin{theorem}[Conditional aggregate closure]
Suppose a uniform exact phase covers every input and every recursively reached
ordinary CNF. At input length $N$, local work is $N^{r+o(1)}$, total executed
child length is at most $N^{v+o(1)}$, and every child has length at most
$N^{h+o(1)}$ for a fixed $0<h<1$. All bounds use one uniform asymptotic
envelope over the same reachable inputs and phases. Fix $N_0$ such that
each child is strictly shorter than its parent for $N>N_0$, and use one
exact terminating base solver on all lengths at most $N_0$.
Repeated child calls count with multiplicity.
Suppose a proved physical simulation incurs at most $N^{g+o(1)}$ cost per
child-machine step, with no additional uncharged semantic call. Then recursive
execution has time $N^{C+o(1)}$, where
\[
 C=\max\left\{1,r,\frac{g+v-h}{1-h}\right\}.
\]
For explicitly materialized children, $g=0$; if local retained state is
$N^{p+o(1)}$ and depth-first composition meets the stated memory model,
peak space is $N^{\max\{p,h\}+o(1)}$.
\end{theorem}
\begin{proof}
For $c\ge1$ and child lengths $m_i$,
\[
 \sum_i m_i^c\le(\max_i m_i)^{c-1}\sum_i m_i
 \le N^{v+h(c-1)+o(1)}.
\]
Physical simulation contributes exponent $g+v+h(c-1)$. Choose any fixed
$c>C$. Then $c>r$ and $c>g+v+h(c-1)$, so standard induction absorbs the
uniform $o(1)$ terms for sufficiently large inputs. Since this holds for
every $c>C$, the exponent bound follows. Along a contracting depth-first
path, retained states and explicit child buffers sum to the stated maximum
exponent. This argument presupposes the actual composable access and memory
implementation; an abstract generator does not establish it.
\end{proof}

\paragraph{Saturation.} At $(r,v,h,p)=(6/5,11/10,1/2,1/2)$, the conditional
pair is $(C,P)=(6/5,1/2)$ when $g=0$. Any conservative per-step access exponent
$g>0$ instead gives $C\ge6/5+2g$. Batch or cache improvements may beat this
simulation bound but need their own proved operation ledger. This is not an
unconditional lower bound on every possible implementation.

\paragraph{Open link CC-A.} Construct the required all-input, successor-closed,
source-visible phase, or the corrected main9 strict successor with
$(r,h,p)=(1,2/3,2/3)$. All acquisition, failed candidates, semantic guards,
encoding, and physical access must fit one simultaneous ledger. Conditional
closure cannot create a transition at a portfolio's UNKNOWN state.

\section{Finite exponent elimination and its limited conclusion}
\begin{lemma}[Fixed-depth elimination]
Suppose one uniform exact SAT-preserving shrinker emits an ordinary CNF of
length $N^{h+o(1)}$, uses time $N^{r+o(1)}$ and space $N^{p+o(1)}$, and has
$0<h<1$. For fixed $d$, apply it $d$ times and then run $A$. Including explicit
successor storage, the resulting bounds are
\[
 T(N)\le N^{\max\{r,h,h^dK\}+o(1)},\qquad
 S(N)\le N^{\max\{p,h,h^d\sigma\}+o(1)}.
\]
\end{lemma}
\begin{proof}
The $j$th successor has length $N^{h^j+o(1)}$. A fixed number of phases has
total cost dominated by the first phase and the final decider call. Retaining
the first explicit child and smaller descendants gives the space maximum.
\end{proof}
For $(r,h,p)=(1,2/3,2/3)$ and finite $K,\sigma$, some fixed $d$ reaches
$(1,2/3)$. Fortnow's near-linear-time/sublinear-space theorem excludes that
pair under the matched random-access model \cite{fortnow2000}. Therefore
the implication from $P=NP$ to existence of the shrinker is itself the
load-bearing unproved statement. The universal product bound
$2\cos(\pi/7)$ is not used.

\section{A separate certificate spine}
Fix one efficiently universal plain-description machine $U$, with polynomial
simulation overhead for each fixed machine. Fix a time bound $t(n)=n^{\omega(1)}$
and a polynomial witness bound $L(n)$. Define a polynomial-time relation
$V(1^n,p,x,\pi)$ and
\[
 I_{n,p}=\{x\in\{0,1\}^n:\exists\pi,\ |\pi|\le L(n),\ V(1^n,p,x,\pi)=1\}.
\]
The same bound $L$ applies throughout. Require, for all $|p|\le n-2$:
\begin{description}
\item[] Completeness: if $U(p)$ halts with the unique final output
$x\in\{0,1\}^n$ within $t(n)$,
then $x\in I_{n,p}$.
\item[] Global ambiguity: $|I_{n,p}|\le2$, after union over all accepted proofs.
\end{description}
Spurious outputs are permitted within the two-element limit. This does not
require exact pointwise endpoint soundness or polynomial-time proof discovery.

\begin{theorem}[Recovered conditional two-list implication]
Such a verifier implies $P\ne NP$.
\end{theorem}
\begin{proof}
For $n\ge2$, let $E_n=\bigcup_{|p|\le n-2}I_{n,p}$. Define
$E=\{x:|x|<2\}\cup\bigcup_{n\ge2}E_n$. The uniform polynomial bound $L$
gives $E\in NP$. There are $2^{n-1}-1$ eligible programs, hence
$|E_n|\le2^n-2$. Completeness implies
\[
 \{0,1\}^*\setminus E\subseteq
 R_{K_U^t}:=\{x:K_U^t(x)\ge |x|-1\}.
\]
Under $P=NP$, the complement of $E$ is an infinite P subset of this set,
contradicting Hirahara's immunity theorem with the same $U,t$ conventions
\cite[Theorem 3.9]{hirahara2020}.
\end{proof}
\paragraph{Open link CC-B.} Existence of $V$ for a suitable fixed $t$ remains
unproved. It is not supplied by hashes, short source code, implicit tableaux,
or tractable-class invariants. The measure is plain time-bounded $K^t$, not
Levin $Kt$. Proof discovery is unnecessary for membership in NP, but must be
charged if a constructive compiler later needs it.

\section{A certificate-length diagonalization boundary}
\begin{theorem}[Exhaustive certificate search obstruction]
Assume a fixed verifier $V$ with polynomially computable witness bound $L(n)$
and verification-time bound $q(n)$ satisfies global ambiguity. There is a
fixed program family $p_n$ of length $O(\log n)$ that outputs an $n$-bit
string outside $E_n$. On a fixed direct machine, its running time is bounded
by
\[
 B(n)=2^{2n+L(n)+O(1)}\operatorname{poly}(n+L(n))\max\{1,q(n)\}.
\]
Let $B_U(n)$ include the efficient universal simulation of this procedure.
For all sufficiently large $n$ with $t(n)\ge B_U(n)$, completeness fails.
Consequently, no polynomial-witness two-list verifier exists for
$t(n)=2^{2^n}$.
\end{theorem}
\begin{proof}
Enumerate $n$-bit strings $x$ lexicographically. For each $x$, enumerate all
$p$ of length at most $n-2$ and all $\pi$ of length at most $L(n)$, checking
$V$. Output the first $x$ with no accepted pair. Global ambiguity guarantees
at least two such strings. At most $2^n(2^{n-1}-1)(2^{L(n)+1}-1)$ verifier
calls suffice. Enumeration overhead yields $B(n)$.
The machine code is fixed and only $n$ needs to be self-delimitingly encoded,
so the resulting $U$ program has length $O(\log n)\le n-2$ for sufficiently
large $n$. If it terminates by $t(n)$, completeness puts its output inside
$E_n$, contradicting its construction. Efficient universal simulation gives
$B_U(n)=2^{\operatorname{poly}(n)}$ for polynomial $L,q$; this is eventually
less than $2^{2^n}$.
\end{proof}
\paragraph{Explicit exponent form.} If $L(n)\le n^a$, $q(n)\le n^b$, and
the fixed universal simulation is bounded by a constant times the $d$th
power of the direct running time, then
\[
 B_U(n)\le 2^{d n^a+2dn+O(\log n)}.
\]
Thus a proposed completeness bound $t(n)=2^{n^c}$ is impossible whenever
$c>\max\{a,1\}$. This is a necessary certificate-degree constraint in that
regime, not a converse existence theorem at equality or below.
\paragraph{Scope.} This refutes an overbroad ``all superpolynomial time
bounds'' certificate claim. It does not refute existence for a slower fixed
bound, such as $n^{\log n}$. A proposed verifier must declare the witness
degree and the decompression time together. The procedure is an elementary
diagonalization, not a SAT lower bound.

\section{Finite hash labels and provenance}
\begin{proposition}[Hash-label capacity]
An exact history identifier over an alphabet of size 62 with $d$ characters
can distinguish at most $62^d$ equivalence classes. Distinguishing $2^k$
classes requires $d\log_2 62\ge k$.
\end{proposition}
\begin{proof} Apply the pigeonhole principle to any injective class identifier.
\end{proof}
The attached 6- and 8-character FNV-1a tools have capacities at most
$35.7252$ and $47.6336$ bits. Extra histories or transcripts may enlarge the
state, but their representation and recovery costs must then be counted.
Collision resistance alone would not imply that a SAT decider must invert a
hash or retain every history.

\paragraph{Exact generator defect.} In the attached randomized generator,
let $s>0$ be cleaned seed length and let $r_0,\ldots,r_{49}$ be the random
pool. At iteration $i=0,\ldots,6$, the index $s+7i$ does not wrap. The $i$
previous deletions shift it to original position $s+8i$, so the output uses
$r_0,r_8,\ldots,r_{48}$ independently of seed contents and positive seed
length. The formatting adds an independent space placement. Thus the seed
does not encode source provenance in this implementation.

\paragraph{Finite experiment.} The companion executable evaluates the actual
HTML scripts. It reproduces identical outputs for distinct nonempty seeds
under identical draws, and an explicit 6-character collision. These finite
observations support the implementation audit, not asymptotic hardness.

\bibliographystyle{plain}
\bibliography{cost_certificate_refs}
\end{document}
